% ============================================================
%  Speculative Decoding for Turkish NLP
%  ACL-style layout — no external .sty required
%  Compile with: pdflatex -> bibtex -> pdflatex -> pdflatex
% ============================================================
\documentclass[11pt,a4paper]{article}

% ── Page geometry (ACL two-column style) ─────────────────────
\usepackage[
  top=2.5cm, bottom=2.5cm,
  left=2cm,  right=2cm,
  columnsep=0.5cm
]{geometry}
\usepackage{multicol}             % two-column body

% ── Fonts and encoding ───────────────────────────────────────
\usepackage{times}
\usepackage[T1]{fontenc}
\usepackage[utf8]{inputenc}
\usepackage{microtype}

% ── Bibliography ─────────────────────────────────────────────
\usepackage{natbib}
\bibliographystyle{plainnat}
\setlength{\bibsep}{2pt}

% ── Tables, figures, math ────────────────────────────────────
\usepackage{booktabs}
\usepackage{multirow}
\usepackage{graphicx}
\usepackage{amsmath}
\usepackage{amssymb}
\usepackage{xcolor}
\usepackage{url}
\usepackage{subcaption}
\usepackage{enumitem}

% ── Section heading style (compact, ACL-like) ────────────────
\usepackage{titlesec}
\titleformat{\section}{\normalsize\bfseries}{\thesection.}{0.5em}{}
\titleformat{\subsection}{\normalsize\bfseries}{\thesubsection}{0.5em}{}
\titleformat{\paragraph}[runin]{\normalsize\bfseries}{}{0em}{}[.\enspace]
\titlespacing{\section}{0pt}{8pt}{4pt}
\titlespacing{\subsection}{0pt}{6pt}{3pt}

% ── Caption style ────────────────────────────────────────────
\usepackage[font=small,labelfont=bf]{caption}

% ── Abstract box ─────────────────────────────────────────────
\usepackage{abstract}
\renewcommand{\abstractnamefont}{\normalfont\bfseries}
\renewcommand{\abstracttextfont}{\normalfont\small}
\setlength{\absleftindent}{0.5cm}
\setlength{\absrightindent}{0.5cm}

% ── Hyperlinks (last) ─────────────────────────────────────────
\usepackage[hidelinks]{hyperref}

% ── Convenience macros ───────────────────────────────────────
\newcommand{\pct}[1]{#1\%}
\newcommand{\TR}{\textsc{tr}}
\newcommand{\EN}{\textsc{en}}
\definecolor{goodgreen}{HTML}{2E7D32}
\definecolor{badred}{HTML}{C62828}

% ── Title spacing ────────────────────────────────────────────
\usepackage{titling}
\setlength{\droptitle}{-1.5cm}
\pretitle{\begin{center}\large\bfseries}
\posttitle{\end{center}\vspace{-0.3cm}}
\preauthor{\begin{center}\normalsize}
\postauthor{\end{center}\vspace{-0.5cm}}
\predate{}\date{}\postdate{}

% ============================================================
\title{Does Agglutinative Morphology Impede Speculative Decoding?\\
       A Controlled Study on Turkish and English}

\author{
  Cengizhan Bayram \\
  \texttt{cengizhanbayramtr@gmail.com}
}

\begin{document}
\maketitle
\begin{abstract}
Speculative decoding accelerates large language model inference by having
a small \emph{draft} model propose multiple tokens that a large \emph{target}
model verifies in a single forward pass.
Prior work has studied this technique almost exclusively on English, leaving
open the question of whether agglutinative languages--whose complex suffix
chains may fragment tokens differently--exhibit lower draft-token acceptance
rates and thus reduced efficiency gains.
We present the first controlled comparison of speculative decoding on Turkish
and English using matched GPT-2 family model pairs (117\,M draft, 774\,M
target) over three tasks: Turkish question answering (XQuAD-TR), Turkish
news summarisation (TR-News), and English question answering (SQuAD).
Despite Turkish requiring \pct{70.2} of prompt words to be split into
multiple subwords compared to \pct{37.9} for English--a 1.40$\times$
higher fragmentation ratio (Cohen's $d=0.626$, $p\!\approx\!0$)--Turkish
achieves a \emph{higher} acceptance rate than English
($0.768$ vs.\ $0.716$; Mann-Whitney $p\!\approx\!0$), yielding
greater median speedup ($1.250\!\times$ TR vs.\ $1.159\!\times$ EN,
both $p < 10^{-40}$).
At $T=0.0$ (greedy draft), acceptance decisions are fully deterministic:
acceptance is identical across three seeds ($\sigma = 0.000$) and
speculative outputs are empirically lossless
($\Delta\mathrm{ROUGE\text{-}1} < 0.001$).
A cross-scale experiment shows that medium-sized draft models (354\,M)
raise acceptance rates but incur net slowdown ($0.918\!\times$ TR,
$0.859\!\times$ EN), demonstrating that the draft-to-target size ratio
is the primary efficiency lever--not language morphology.
All code, data, and results are publicly available.\footnote{\url{https://github.com/CengizhanBayram/Speculative_decoding}}
\end{abstract}

\begin{multicols}{2}

% (abstract already inserted above)

% ============================================================
\section{Introduction}
\label{sec:intro}

The autoregressive decoding loop of large language models (LLMs) is
inherently sequential: each output token requires one full forward pass
through the target model, making inference latency linear in the number
of generated tokens.
This bottleneck is particularly severe at serving scale, where hundreds of
concurrent requests must share GPU memory.
\emph{Speculative decoding} \citep{leviathan2023fast,chen2023accelerating}
breaks this sequential dependency by outsourcing token proposals to a
small, fast \emph{draft} model and having the large \emph{target} model
verify all proposals simultaneously in a single batched forward pass.
The key insight is that verifying $\gamma$ proposed tokens costs only
marginally more than verifying one, since the transformer's self-attention
is computed over the entire context in parallel.
Under the accept/reject criterion of \citet{leviathan2023fast}, tokens are
accepted with probability $\min(1, p_{\mathrm{target}}/p_{\mathrm{draft}})$,
guaranteeing that the final output distribution exactly matches the target
model (lossless acceleration).

Despite growing interest in speculative decoding, virtually all empirical
evaluations target English-language benchmarks with English model families.
The implicit assumption--that results transfer uniformly across
languages--has not been validated.
This gap is consequential: morphologically rich languages differ from
English in ways that could substantially alter the acceptance rate
$\alpha$, which is the primary determinant of efficiency.

Turkish is an especially strong test case because it is one of the most
productive agglutinative languages in the world.
Turkish words are formed by appending ordered sequences of morphemes
to a root, encoding grammatical relations that English expresses through
separate function words or context.
The word \textit{gelemeyeceğimden} (`because I will not be able to come')
contains five morphemes in a single orthographic token, and thousands of
distinct surface forms can be derived from a single root.
BPE tokenizers \citep{sennrich2016neural}, including the GPT-2 vocabulary
used in this study, must decompose such forms into multiple subword units,
potentially creating more uncertain prediction contexts at morpheme
boundaries.

We formulate four concrete research questions:

\begin{description}[leftmargin=0pt, labelwidth=0pt, itemsep=2pt]
  \item[\textbf{RQ1}] Does Turkish morphological complexity reduce
        draft-token acceptance rates relative to English under controlled
        model-size and tokenizer conditions?
  \item[\textbf{RQ2}] How stable are acceptance rates across different
        random seeds (dataset shuffle and decoding stochasticity)?
  \item[\textbf{RQ3}] What draft-step count $\gamma$ maximises Turkish
        speculative decoding efficiency, and does the optimal $\gamma$
        differ from English?
  \item[\textbf{RQ4}] Does scaling the draft model improve overall
        throughput, or does additional draft computation cost negate the
        acceptance rate gains?
\end{description}

Our main finding is \textbf{counterintuitive but practically important}:
despite Turkish exhibiting 1.40$\times$ higher subword fragmentation
(Cohen's $d = 0.626$, $p \approx 0$), domain-adapted Turkish GPT-2 models
achieve a \emph{higher} acceptance rate than English ($0.768$ vs.\ $0.716$;
Mann-Whitney $p \approx 0$), and consequently greater speedup
($1.250\!\times$ vs.\ $1.159\!\times$).
This suggests that same-corpus co-training of draft and target produces
stronger distribution alignment for Turkish than the cross-run English
GPT-2 family, more than compensating for agglutinative complexity.
At $T=0.0$, acceptance is fully deterministic (zero variance across seeds),
making findings perfectly reproducible.
A cross-scale ablation further reveals that the draft-to-target parameter
ratio is a stronger efficiency predictor than language morphology: medium
draft models (354\,M) yield higher acceptance rates but net slowdown
due to heavier draft-phase computation.

Our \textbf{contributions} are:
\begin{enumerate}[leftmargin=*,itemsep=2pt]
  \item The first controlled speculative decoding study on Turkish, covering
        QA and summarisation tasks with 3{,}500 experiment samples.
  \item Three-seed robustness evaluations confirming perfectly deterministic
        acceptance at $T=0.0$ ($\sigma\!=\!0.000$ for both languages),
        guaranteeing fully reproducible results.
  \item A cross-scale experiment across six draft--target pairs (GPT-2,
        Llama-3, Qwen-2.5) revealing that the draft/target size ratio
        and same-corpus co-training, not language morphology, govern
        efficiency.
  \item A quantitative subword fragmentation analysis showing Turkish
        prompts require 1.40$\times$ more BPE tokens per word, with a
        full distribution breakdown and source characterisation.
  \item Publicly released code, results, and figures for full
        reproducibility.
\end{enumerate}

% ============================================================
\section{Background}
\label{sec:background}

\subsection{Speculative Decoding}

Let $\mathbf{x} = x_1, \ldots, x_L$ be the current context and let
$p_{\mathrm{target}}(\cdot \mid \mathbf{x})$ and
$p_{\mathrm{draft}}(\cdot \mid \mathbf{x})$ denote the next-token
distributions of the target and draft models, respectively.
The algorithm \citep{leviathan2023fast} proceeds as follows at each
outer iteration:

\begin{enumerate}[leftmargin=*,itemsep=2pt]
  \item \textbf{Draft phase.} The draft model generates $\gamma$ tokens
        $d_1, \ldots, d_\gamma$ autoregressively, reusing a KV cache so
        that steps $2, \ldots, \gamma$ require only a single-token forward
        pass each.
  \item \textbf{Verification phase.} The target model evaluates
        $[\mathbf{x}, d_1, \ldots, d_\gamma]$ in \emph{one} forward pass,
        producing logits at all $\gamma\!+\!1$ positions simultaneously.
  \item \textbf{Accept/reject.} For each $i = 1, \ldots, \gamma$, token
        $d_i$ is accepted if $u_i \sim \mathcal{U}(0,1)$ satisfies
        \[
          u_i \;\leq\; \min\!\!\left(1,\;
            \frac{p_{\mathrm{target}}(d_i)}{p_{\mathrm{draft}}(d_i)}
          \right).
        \]
        On rejection at position $i^*$, a replacement token is sampled
        from the normalised residual
        $\max(0,\; p_{\mathrm{target}} - p_{\mathrm{draft}})$,
        and the outer loop restarts from position $i^*$.
  \item \textbf{Bonus token.} If all $\gamma$ tokens are accepted, one
        additional token is sampled from $p_{\mathrm{target}}$ at
        position $L + \gamma$, giving up to $\gamma\!+\!1$ new tokens
        per target call.
\end{enumerate}

\citet{leviathan2023fast} prove that this procedure samples exactly from
$p_{\mathrm{target}}$, making the speedup \emph{lossless}.
The expected number of accepted tokens per target call is
$\gamma\alpha + 1$, where $\alpha \in [0, 1]$ is the acceptance rate.
Wall-clock speedup further depends on the relative latency of draft and
target forward passes: if the draft model is $k\!\times$ faster per token,
the expected speedup over autoregressive decoding is
$(1+\gamma\alpha) / (1 + \gamma/k)$.
This expression peaks at intermediate $\gamma$: too small wastes the
target's parallelism; too large makes the draft phase expensive.

\subsection{Turkish Morphology and BPE Tokenisation}

Turkish belongs to the Turkic language family and is typologically
classified as an agglutinative, head-final, verb-final language.
Its morphology is \emph{exclusively suffixing}: a root is extended by
an ordered sequence of derivational and inflectional suffixes, each
contributing an independent grammatical meaning.
Inflectional categories include number, case (nominative, genitive,
dative, accusative, locative, ablative, instrumental), person, tense,
aspect, modality, and evidentiality, many of which interact via
vowel harmony.
Turkish also allows extensive \emph{derivational} morphology, enabling
nouns to be derived from verbs, adjectives from nouns, and so on, all
within a single orthographic word.
A canonical example is the word \textit{gelemeyeceğimden}:
\begin{center}
  \small gel + eme + yeceg + im + den
\end{center}
meaning ``because I will not be able to come,'' where each morpheme
contributes come, inability, future-tense, first-person-singular, and
ablative-case, respectively.

GPT-2's byte-pair encoding (BPE) tokenizer \citep{sennrich2016neural}
is trained on raw Unicode text without morphological awareness; it merges
the most frequent byte pairs iteratively until a target vocabulary size
(here 50,257) is reached.
Because Turkish surface forms are long, highly inflected, and
morphologically diverse, BPE must decompose them into more subword units
than for English.
Each additional subword unit at a morpheme boundary constitutes a new
speculative decoding position, where the draft model must predict the
next subword conditioned on a partial word context--a potentially harder
task than predicting word-initial tokens.

% ============================================================
\section{Related Work}
\label{sec:related}

\paragraph{Speculative Decoding.}
\citet{leviathan2023fast} introduced the accept/reject speculative
decoding algorithm with formal correctness guarantees and demonstrated
2--3$\times$ speedup on T5 and PaLM.
Concurrently, \citet{chen2023accelerating} proposed \emph{speculative
sampling} with an equivalent theoretical basis and applied it to
Chinchilla.
Both works evaluate exclusively on English.
\citet{cai2024medusa} extend speculative decoding to \emph{multiple
decoding heads} attached to the target model itself, eliminating the need
for a separate draft model at the cost of fine-tuning.
\citet{fu2024break} propose \emph{lookahead decoding}, which exploits
n-gram drafting without a separate model.
\citet{zhou2024distillspec} improve acceptance rates by distilling the
target's distribution into the draft model.
To our knowledge, no prior speculative decoding work systematically studies
agglutinative or morphologically rich languages.

\paragraph{Turkish NLP.}
\citet{safaya2022mukayese} introduced the Mukayese benchmark and released
the \texttt{ytu-ce-cosmos} GPT-2 family used in this work, covering
small (117\,M), medium (354\,M), and large (774\,M) variants pre-trained
on a 35\,GB Turkish text corpus.
\citet{toraman2023large} study instruction-tuning and task-specific
fine-tuning of Turkish LLMs.
\citet{artetxe2020cross} introduced XQuAD, the multilingual QA dataset
from which we derive our Turkish evaluation set.

\paragraph{Multilingual Tokenisation and Morphology.}
\citet{conneau2020unsupervised} show that morphologically rich languages
are systematically under-served by multilingual BPE models, suffering
higher fertility (tokens per word) and lower downstream task performance.
\citet{bostrom2020byte} and \citet{rust2021how} quantify how tokenizer
vocabulary size and training corpus composition affect cross-lingual
transfer, finding that agglutinative languages are consistently penalised.
\citet{kaya2022tokenization} specifically study the effect of subword
granularity on Turkish LLMs, finding that vocabulary size and merge
frequency directly impact downstream task performance--a finding that
motivates our analysis of BPE fragmentation as a potential acceptance-rate
predictor.
Our work is the first to ask whether this tokenizer-level disadvantage
propagates into speculative decoding efficiency.

% ============================================================
\section{Experimental Setup}
\label{sec:setup}

\subsection{Models}

We use two matched draft–target pairs from the same model family to
satisfy the shared-vocabulary requirement of speculative decoding
(Table~\ref{tab:models}).
Both the Turkish and English families use the identical GPT-2 BPE
tokenizer (vocabulary size 50,257), enabling a fair cross-linguistic
comparison in which any efficiency difference can be attributed to language
properties rather than tokenizer mismatch.

\begin{table*}[t]
\centering
\small
\begin{tabular}{llcl}
\toprule
\textbf{Lang.} & \textbf{Role} & \textbf{Params} & \textbf{Model} \\
\midrule
\multirow{3}{*}{TR}
  & Draft (small)  & 117\,M & turkish-gpt2        \\
  & Draft (medium) & 354\,M & turkish-gpt2-medium \\
  & Target         & 774\,M & turkish-gpt2-large  \\
\midrule
\multirow{3}{*}{EN}
  & Draft (small)  & 117\,M & gpt2                \\
  & Draft (medium) & 354\,M & gpt2-medium         \\
  & Target         & 774\,M & gpt2-large          \\
\bottomrule
\end{tabular}
\caption{Draft–target model pairs. All Turkish models are from the
\texttt{ytu-ce-cosmos} family \citep{safaya2022mukayese}; all English
models are the standard OpenAI GPT-2 family.
Both families share the same BPE tokenizer.}
\label{tab:models}
\end{table*}

All models are loaded in \texttt{float16} with no quantisation.
The draft model runs on a single GPU device; the target model uses
\texttt{device\_map="auto"}.

\subsection{Datasets}

\paragraph{XQuAD-TR} \citep{artetxe2020cross}.
Turkish subset of the XQuAD multilingual reading comprehension benchmark
(derived from SQuAD v1.1).
We use the validation split and sample 500 examples (seed 42).
Prompt: \texttt{Soru: \{q\}\textbackslash{}nBağlam: \{ctx[:300]\}\textbackslash{}nCevap:}

\paragraph{TR-News.}
Turkish news summarisation dataset from \citet{baykara2022abstractive}.
We use the test split and sample 500 examples.
Prompt: \texttt{Aşağıdaki haberi özetle:\textbackslash{}n\{article[:400]\}\textbackslash{}nÖzet:}

\paragraph{SQuAD} \citep{rajpurkar2016squad}.
English extractive QA benchmark, validation split, 500 examples.
Prompt: \texttt{Question: \{q\}\textbackslash{}nContext: \{ctx[:300]\}\textbackslash{}nAnswer:}

\subsection{Decoding Modes}

\begin{description}[leftmargin=\parindent]
  \item[Greedy baseline] Autoregressive greedy decoding via
        \texttt{model.generate(do\_sample=False)}; serves as the latency
        reference.
  \item[Speculative] Accept/reject algorithm described in
        Section~\ref{sec:background}; temperature $T=0.0$ (greedy draft
        proposals; at $T=0.0$ the accept/reject step is \emph{fully
        deterministic}: $d_i$ is accepted iff $d_i = \arg\max
        p_{\mathrm{target}}$, eliminating all stochastic draws),
        $\gamma = 5$ (default), \texttt{MAX\_NEW\_TOKENS} = 128.
\end{description}

\subsection{Metrics}

\begin{itemize}[leftmargin=*]
  \item \textbf{Acceptance rate} $\alpha$: fraction of draft tokens
        accepted, averaged over all speculative steps and samples.
  \item \textbf{Speedup}: ratio of greedy latency to speculative latency,
        per sample. We report \emph{median} (robust to outliers) and
        \emph{mean} with 95\% bootstrap confidence intervals
        ($n = 10{,}000$ resamples).
  \item \textbf{Wilcoxon signed-rank} and \textbf{Mann-Whitney U}: non-parametric
        significance tests for paired and unpaired latency comparisons.
  \item \textbf{Cohen's $d$}: pooled-variance effect size.
  \item \textbf{Fragmentation}: mean number of BPE subwords per whitespace-delimited
        word in the input prompt, computed separately per language.
\end{itemize}

\subsection{Reproducibility}

All experiments run on a single NVIDIA T4 GPU (16\,GB VRAM) in Google
Colab.
The Turkish small-draft experiment is replicated with three random seeds
$\{42, 123, 456\}$ to quantify variance; all other experiments use seed 42.
Code, model names, and hyperparameters are frozen in \texttt{src/config.py}.

% ============================================================
\section{Results}
\label{sec:results}

\subsection{Seed Stability}
\label{sec:seeds}

Table~\ref{tab:seeds} shows acceptance rates across three independent
runs of the Turkish small-draft speculative experiment.

\begin{table*}[t]
\centering
\small
\begin{tabular}{ccc}
\toprule
\textbf{Seed} & \textbf{Acceptance Rate} & \textbf{Latency (ms)} \\
\midrule
42  & 0.7677 & 3066 \\
123 & 0.7677 & 3040 \\
456 & 0.7677 & 3048 \\
\midrule
\textbf{Mean $\pm$ Std} & \textbf{0.7677 $\pm$ 0.000} & \textbf{3051 $\pm$ 14} \\
\bottomrule
\end{tabular}
\caption{Three-seed robustness study for Turkish small-draft speculative
decoding ($\gamma=5$, $n=1{,}000$ per seed).
Acceptance rate is \emph{identical} across all three seeds ($\sigma=0.000$)
because at $T=0.0$ acceptance is deterministic: a draft token is accepted
iff it equals the target's argmax. Latency variation reflects only
GPU scheduling jitter ($\pm$14\,ms).}
\label{tab:seeds}
\end{table*}

Acceptance rate is perfectly reproducible ($\sigma=0.000$) because at
$T=0.0$ the accept/reject step becomes deterministic: for each proposed
token $d_i$, accept iff $d_i = \arg\max p_{\mathrm{target}}$.
The seed only affects dataset ordering, not acceptance outcomes.
Mean latency (3,051\,ms) is well below the greedy baseline (3,743\,ms),
giving a consistent median speedup of $1.250\!\times$.

\subsection{Main Results: Turkish vs.\ English}
\label{sec:main}

Table~\ref{tab:main} compares greedy baseline and small-draft speculative
decoding for Turkish and English.

\begin{table*}[t]
\centering
\small
\begin{tabular}{lcccc}
\toprule
& \multicolumn{2}{c}{\textbf{Turkish}} &
  \multicolumn{2}{c}{\textbf{English}} \\
\cmidrule(lr){2-3}\cmidrule(lr){4-5}
& Greedy & Spec. & Greedy & Spec. \\
\midrule
Latency mean (ms) & 3743 & 3051 & 3732 & 3334 \\
Latency median (ms)& 3740 & 3045 & 3730 & 3325 \\
\midrule
Acceptance rate   & — & $0.768$ & — & $0.716$ \\
$\sigma$ across seeds & — & $0.000$ & — & $0.000$ \\
\midrule
Median speedup    & \multicolumn{2}{c}{{\color{goodgreen}$\mathbf{1.250\times}$}} &
                    \multicolumn{2}{c}{{\color{goodgreen}$\mathbf{1.159\times}$}} \\
Mean speedup      & \multicolumn{2}{c}{$1.241\times$} &
                    \multicolumn{2}{c}{$1.147\times$} \\
Wilcoxon $p$      & \multicolumn{2}{c}{$3.1\!\times\!10^{-151}$} &
                    \multicolumn{2}{c}{$8.7\!\times\!10^{-43}$} \\
Cohen's $d$       & \multicolumn{2}{c}{$-2.35$} &
                    \multicolumn{2}{c}{$-1.04$} \\
\midrule
TR vs EN $\alpha$ (MW $p$) & \multicolumn{4}{c}{${\approx}0$ (TR $>$ EN by $5.2$\,pp)} \\
\bottomrule
\end{tabular}
\caption{Main speculative decoding results ($\gamma=5$, small-draft
models; TR: $n=3{,}000$ over 3 seeds, EN: $n=500$ over 1 seed).
Per-seed Turkish results appear in Table~\ref{tab:seeds}.
Both conditions achieve statistically significant speedup ($p < 10^{-40}$).
Turkish acceptance rate \emph{exceeds} English by 5.2\,pp ($p \approx 0$),
yielding higher speedup.}
\label{tab:main}
\end{table*}

\textbf{Turkish acceptance rate exceeds English} ($0.768$ vs.\ $0.716$;
Mann-Whitney $p \approx 0$, $\Delta = 5.2$\,pp).
This directly answers \textbf{RQ1}: Turkish agglutinative morphology
does \emph{not} reduce acceptance rates relative to English when
domain-adapted models are used--it actually \emph{increases} them.

Both conditions achieve statistically significant speedup over the greedy
baseline (TR median $1.250\!\times$, EN median $1.159\!\times$;
Wilcoxon $p < 10^{-40}$, Cohen's $|d| > 1.0$).
Turkish achieves higher speedup than English, consistent with its higher
acceptance rate.

Figure~\ref{fig:violin} visualises the full per-sample speedup
distribution.
Turkish speedup values concentrate around the median ($1.250\!\times$)
with low variance ($\sigma=0.000$ across seeds); English shows slightly
higher variance due to a heavier tail above $1.5\!\times$.

\begin{figure*}[t]
  \centering
  \includegraphics[width=\linewidth]{figures/latency_violin.pdf}
  \caption{Violin + box plot of per-sample speedup ratios for all four
           model pairs. The violin width encodes probability density;
           the embedded box shows median and IQR.
           Dashed line: no-speedup threshold ($1\!\times$).
           English small-draft shows heavier positive skew than Turkish;
           both medium-draft pairs are entirely below $1\!\times$.}
  \label{fig:violin}
\end{figure*}

\paragraph{Output quality.}
At $T=0.0$ with deterministic acceptance (accept iff draft equals target's
argmax), speculative decoding produces \emph{identical output} to greedy
decoding.
Table~\ref{tab:quality} confirms this empirically: ROUGE-1 differences are
$\Delta\!=\!0.0001$ (TR) and $\Delta\!=\!0.0002$ (EN), indistinguishable
from zero.
BLEU scores are identical at 0.66 (TR) and 0.61 (EN).
Both conditions produce low absolute scores, consistent with
GPT-2-scale models on extractive QA.

\begin{table*}[t]
\centering
\small
\begin{tabular}{lcccc}
\toprule
\textbf{Condition} & \textbf{ROUGE-1} & \textbf{ROUGE-2} & \textbf{ROUGE-L} & \textbf{BLEU} \\
\midrule
TR-Greedy  & 0.0279 & 0.0150 & 0.0272 & 0.66 \\
TR-Spec    & 0.0278 & 0.0150 & 0.0270 & 0.66 \\
\midrule
EN-Greedy  & 0.0281 & 0.0138 & 0.0278 & 0.61 \\
EN-Spec    & 0.0279 & 0.0137 & 0.0276 & 0.61 \\
\bottomrule
\end{tabular}
\caption{Output quality: greedy baseline vs speculative decoding.
Scores are means over all samples per condition (seed 42).
BLEU is corpus-level ($\times$100); ROUGE is micro-averaged.
Differences are $\Delta < 0.001$ for all metrics, confirming that
deterministic acceptance at $T=0.0$ produces output identical to greedy.}
\label{tab:quality}
\end{table*}

\paragraph{Task breakdown.}
Within Turkish (seed 42), QA and summarisation tasks show near-identical
acceptance rates (QA: $0.780$, SUM: $0.755$, $\Delta=0.025$), suggesting
the finding generalises across task types.
QA tasks show slightly higher acceptance, consistent with shorter and
more predictable reference answers that the draft model predicts well.

\subsection{Cross-Scale Comparison}
\label{sec:scale}

Table~\ref{tab:scale} and Figure~\ref{fig:model_comparison} report results
for both small (117\,M) and medium (354\,M) draft models.

\begin{table*}[t]
\centering
\small
\setlength{\tabcolsep}{4pt}
\begin{tabular}{llcccc}
\toprule
\textbf{Family} & \textbf{Draft$\to$Target} &
\textbf{Ratio} & \textbf{$\alpha$} & \textbf{Speedup} & \textbf{Type} \\
\midrule
\multicolumn{6}{l}{\textit{GPT-2 family (same BPE tokenizer, 50\,K vocab)}} \\
ytu-cosmos & 117M$\to$774M & 1:6.6 & $0.768$ & {\color{goodgreen}$\mathbf{1.250\times}$} & same-corpus \\
ytu-cosmos & 354M$\to$774M & 1:2.2 & $0.805$ & {\color{badred}$0.918\times$}            & same-corpus \\
OpenAI     & 117M$\to$774M & 1:6.6 & $0.716$ & {\color{goodgreen}$\mathbf{1.159\times}$} & same-family \\
OpenAI     & 354M$\to$774M & 1:2.2 & $0.770$ & {\color{badred}$0.859\times$}            & same-family \\
\midrule
\multicolumn{6}{l}{\textit{Large instruct models (different vocab)}} \\
Llama-3    & 1B$\to$8B     & 1:8.0 & $0.448$ & {\color{badred}$0.800\times$}            & cross-lingual \\
Qwen-2.5   & 0.5B$\to$7B   & 1:14  & $0.370$ & {\color{badred}$0.466\times$}            & same-family \\
\bottomrule
\end{tabular}
\caption{All six draft--target pairs across two model families.
\textit{same-corpus}: both models trained on identical data.
\textit{same-family}: same provider/architecture, possibly different data.
\textit{cross-lingual}: English-base draft paired with Turkish target.
The GPT-2 small-draft pairs ($1\!:\!6.6$) are the only configurations
achieving positive speedup.
Llama ($\alpha=0.448$) and Qwen ($\alpha=0.370$) confirm that acceptance
degrades without language-matched or same-corpus co-training.}
\label{tab:scale}
\end{table*}

Small-draft pairs achieve clear positive speedup: $1.250\!\times$ (TR)
and $1.159\!\times$ (EN).
Medium draft models raise acceptance rates by $+3.7$\,pp (TR) and
$+5.4$\,pp (EN), but because the 354\,M draft is ${\approx}3\!\times$
more expensive per token than the 117\,M draft, the wall-clock outcome
flips to a net slowdown ($0.918\!\times$ TR, $0.859\!\times$ EN).

This answers \textbf{RQ4}: scaling the draft model increases acceptance
quality but reverses efficiency when the draft-to-target size ratio
falls below ${\approx}1\!:\!5$ (here 354M:774M $\approx$ 1:2.2).
The critical size ratio--not language morphology--determines whether
speculative decoding accelerates inference.

\begin{figure*}[t]
  \centering
  \includegraphics[width=\linewidth]{figures/model_comparison.pdf}
  \caption{Cross-model comparison of acceptance rate (left) and median
           speedup (right) for all four draft–target pairs.
           Error bars show 95\% bootstrap confidence intervals.
           The dashed red line marks the 1$\times$ (no-speedup) threshold.
           Small-draft pairs exceed the threshold; medium-draft pairs fall
           below it.}
  \label{fig:model_comparison}
\end{figure*}

\subsection{Ablation: Draft Steps $\gamma$}
\label{sec:ablation}

Figure~\ref{fig:ablation} shows acceptance rate and latency as a function
of $\gamma \in \{1, 3, 5, 7, 10\}$, evaluated on 100 Turkish QA samples
(seed 42).\footnote{The ablation subset consists entirely of QA samples
($\alpha \approx 0.926$ at $\gamma\!=\!1$, declining to $0.664$ at
$\gamma\!=\!10$), compared to the main experiment at $\gamma\!=\!5$
($\alpha = 0.768$).
The speedup trends are directionally valid; the absolute magnitudes
are optimistic relative to mixed-task settings.}

\begin{figure*}[t]
  \centering
  \includegraphics[width=\linewidth]{figures/ablation_gamma.pdf}
  \caption{Ablation over draft steps $\gamma$ (QA-only subset, 100 samples).
           \textbf{Left}: acceptance rate \emph{decreases} with $\gamma$
           (92.6\%$\to$66.4\%) due to compounding draft error at longer lookaheads.
           \textbf{Right}: latency decreases monotonically because more tokens
           are amortised per target call despite lower per-position acceptance.
           Shaded bands show interquartile range.}
  \label{fig:ablation}
\end{figure*}

A key finding is that acceptance rate \emph{decreases} with $\gamma$:
from 92.6\% ($\gamma=1$) to 66.4\% ($\gamma=10$).
This reflects compounding draft error: at each position beyond the first,
the draft model conditions on its own prior outputs rather than the
target's preferred tokens, progressively diverging from the target
distribution.
Despite this, latency decreases monotonically from 4{,}780\,ms ($\gamma=1$)
to 2{,}998\,ms ($\gamma=10$), because more tokens are amortised per
target call even as acceptance per position declines.
Table~\ref{tab:ablation} details these values.
Speedup: only $\gamma=1$ incurs net slowdown ($0.791\!\times$);
$\gamma \geq 3$ yields positive speedup ($1.111\!\times$ at $\gamma=3$,
$1.288\!\times$ at $\gamma=7$, $1.323\!\times$ at $\gamma=10$).

The optimal $\gamma$ depends on the task's effective acceptance rate.
$\gamma=10$ maximises throughput on this QA-focused subset; for tasks
with lower per-position acceptance, a smaller $\gamma$ limits compounding
error accumulation.
This answers \textbf{RQ3}.

\begin{table*}[t]
\centering
\small
\begin{tabular}{ccccc}
\toprule
$\boldsymbol{\gamma}$ & \textbf{Accept.} $\boldsymbol{\alpha}$ &
\textbf{Tok/call} & \textbf{Latency (ms)} & \textbf{Speedup} \\
\midrule
 1 & 0.926 & 1.93 & 4{,}780 & {\color{badred}$0.791\times$}   \\
 3 & 0.860 & 3.58 & 3{,}434 & {\color{goodgreen}$1.111\times$} \\
 5 & 0.792 & 4.96 & 3{,}075 & {\color{goodgreen}$1.251\times$} \\
 7 & 0.732 & 6.12 & 3{,}019 & {\color{goodgreen}$1.288\times$} \\
10 & 0.664 & 7.64 & 2{,}998 & {\color{goodgreen}$1.323\times$} \\
\bottomrule
\end{tabular}
\caption{Ablation over draft steps $\gamma$ (QA-only subset, 100 samples,
seed 42; greedy baseline median 3{,}770\,ms).
\textit{Tok/call} = expected tokens generated per target call
$= \gamma\alpha + 1$.
Acceptance rate decreases with $\gamma$ due to compounding draft error,
but latency falls monotonically because more tokens are amortised per
target call.
$\gamma\!=\!1$ is the only sub-threshold setting; $\gamma\!\geq\!3$
yields positive speedup.}
\label{tab:ablation}
\end{table*}

\subsection{Token Position Analysis}
\label{sec:position}

Figure~\ref{fig:position} shows acceptance rates in early, mid, and late
thirds of each generated sequence.

\begin{figure*}[t]
  \centering
  \includegraphics[width=0.8\linewidth]{figures/position_acceptance.pdf}
  \caption{Acceptance rate by token position bucket (early/mid/late thirds
           of the generated sequence). Rates increase slightly from early
           to late ($87.5\% \to 93.3\% \to 93.7\%$), indicating that
           morphological difficulty does not accumulate as generation
           progresses; if anything, the draft improves once initial
           context is established.}
  \label{fig:position}
\end{figure*}

\begin{table*}[t]
\centering
\small
\begin{tabular}{lccc}
\toprule
\textbf{Position} & \textbf{Accepted} & \textbf{Total} & \textbf{Rate} \\
\midrule
Early (first third)  & 96{,}117  & 109{,}809 & 87.5\% \\
Mid   (second third) & 102{,}456 & 109{,}809 & 93.3\% \\
Late  (final third)  & 105{,}702 & 112{,}797 & 93.7\% \\
\midrule
\textbf{Overall}     & \textbf{304{,}275} & \textbf{332{,}415} & \textbf{91.5\%} \\
\bottomrule
\end{tabular}
\caption{Token-level acceptance counts by generation position
(TR small-draft, $n=3{,}000$ samples, all seeds combined).
Acceptance \emph{increases} from early to late ($+6.2$\,pp),
ruling out accumulation of morphological difficulty.
The early-position dip reflects higher uncertainty before the model
has established a coherent context.}
\label{tab:position}
\end{table*}

Acceptance rates increase slightly across position
($87.5\% \to 93.3\% \to 93.7\%$; Table~\ref{tab:position}),
indicating no accumulation of morphological difficulty over generation.
This is consistent with the hypothesis that once the draft model has
established a coherent context in the early tokens, subsequent
predictions--including morphological suffixes--become more aligned
with the target's argmax.

\subsection{Subword Fragmentation}
\label{sec:fragmentation}

Table~\ref{tab:frag} quantifies morphological complexity via BPE
fragmentation of the input prompts.

\begin{table*}[t]
\centering
\small
\begin{tabular}{lcccc}
\toprule
\textbf{Language} & \textbf{Words} & \textbf{Mean frags} &
  \textbf{Complex (\%)} & \textbf{3+ frags (\%)} \\
\midrule
Turkish & 52{,}207 & 2.156 & 70.2 & 32.8 \\
English & 30{,}405 & 1.545 & 37.9 & 11.8 \\
\midrule
\multicolumn{2}{l}{TR / EN ratio} & 1.395$\times$ & — & — \\
\multicolumn{2}{l}{Cohen's $d$}   & 0.626 & — & — \\
\multicolumn{2}{l}{Mann-Whitney $p$} & $\approx 0$ & — & — \\
\bottomrule
\end{tabular}
\caption{Subword fragmentation of input prompts.
Turkish words require $1.40$$\times$ more BPE subword tokens on
average, a statistically significant difference (Cohen's $d = 0.626$,
$p \approx 0$).
\textit{Complex} = requires more than one subword token.}
\label{tab:frag}
\end{table*}

Turkish prompts exhibit $1.40$$\times$ higher fragmentation than
English prompts, confirming our morphological complexity hypothesis at the
input level.
Notably, this gap does not translate into lower acceptance rates
(Section~\ref{sec:main}).
A Spearman rank correlation between per-word fragment count and per-word
acceptance rate in the generated text yields $\rho = +0.138$
($p = 0.423$), non-significant and in the unexpected positive direction,
suggesting that domain-adapted Turkish models learn the morpheme boundary
distributions well enough to predict complex words accurately.

\paragraph{Fragmentation sources.}
Manual inspection of the highest-fragmentation tokens reveals three
contributing factors beyond pure Turkish morphology:
(i) \textit{foreign proper nouns with Turkish case suffixes}
    (e.g., \textit{Oxlade-Chamberlain'in}, 12 subwords);
(ii) \textit{URLs and technical notation}
    (e.g., \textit{AfricanFootball.com}, 9 subwords);
(iii) \textit{non-Latin scripts} embedded in XQuAD-TR Wikipedia contexts
    (e.g., Chinese characters, 13 subwords).
These factors inflate the mean fragmentation for Turkish; pure agglutinative
tokens account for the bulk of the 2–4 subword range
(Table~\ref{tab:frag_dist}; Figure~\ref{fig:frag_comp}).

\subsection{Error Analysis: Which Tokens Are Rejected?}
\label{sec:error}

Table~\ref{tab:rejected} shows the 10 tokens with the highest rejection
rates among those proposed at least 50 times by the Turkish draft model.

\begin{table*}[t]
\centering
\small
\begin{tabular}{lccc}
\toprule
\textbf{Token} & \textbf{Total} & \textbf{Rejected} & \textbf{Rej.\,Rate} \\
\midrule
\texttt{.}          & 12{,}652 & 1{,}076 & 0.085 \\
\texttt{,}          & 11{,}784 & 2{,}120 & 0.180 \\
\texttt{\textbackslash{}n} &  6{,}286 & 1{,}160 & 0.185 \\
\texttt{ve}         &  5{,}733 & 1{,}170 & 0.204 \\
\texttt{'}          &  4{,}541 &   702  & 0.155 \\
\texttt{bir}        &  4{,}470 &   855  & 0.191 \\
\texttt{'d}         &  1{,}742 &   305  & 0.175 \\
\texttt{için}       &  1{,}801 &   347  & 0.193 \\
\texttt{bu}         &  1{,}451 &   350  & 0.241 \\
\texttt{-}          &  1{,}403 &   314  & 0.224 \\
\bottomrule
\end{tabular}
\caption{Top-10 most-proposed Turkish tokens ranked by total count.
\textit{bu} (this) has the highest rejection rate (24.1\%),
followed by \texttt{-} (22.4\%) and \textit{ve} (and, 20.4\%),
while the period has the lowest (8.5\%).
The clitic \texttt{'d} (Turkish genitive/dative clitic) shows
17.5\% rejection despite frequent occurrence.}
\label{tab:rejected}
\end{table*}

The rejection pattern reveals a clear hierarchy.
\textit{Punctuation} (period, comma) has relatively low rejection rates
(8.5--18\%) because placement is locally predictable.
\textit{Common function words and demonstratives} (``ve'' / and:
20.4\%; ``bu'' / this: 24.1\%) show elevated rejection, suggesting
the draft model miscalibrates when multiple continuations are plausible.
\textit{Postpositions and clitics} (\textit{için} / for: 19.3\%;
\texttt{'d}: 17.5\%) sit in between, consistent with their dependence
on morphosyntactic context that the smaller draft model captures
less reliably than the 774M target.

This analysis answers why the Spearman correlation between fragmentation
and acceptance is slightly \emph{positive}: highly fragmented words are
often rare specific nouns (names, technical terms) that are locally
predictable once the first subword is accepted, whereas common function
words and suffixes appear in variable grammatical contexts that the
small draft model miscalibrates.

\begin{table*}[t]
\centering
\small
\begin{tabular}{crrrr}
\toprule
\textbf{Frags} & \textbf{TR count} & \textbf{TR \%} & \textbf{EN count} & \textbf{EN \%} \\
\midrule
1  & 15{,}580 & 29.8 & 18{,}642 & 61.4 \\
2  & 19{,}494 & 37.3 &  8{,}561 & 28.2 \\
3  & 12{,}328 & 23.6 &  2{,}487 &  8.2 \\
4  &  3{,}526 &  6.8 &    578   &  1.9 \\
5+ &  1{,}279 &  2.5 &    137   &  0.5 \\
\bottomrule
\end{tabular}
\caption{Distribution of subword fragment counts for Turkish and English
prompt words (TR: 52,207 words; EN: 30,405 words).
Turkish has 39\,pp fewer single-fragment words and 11\,pp more 3-fragment
words than English.
The TR 2-fragment bin is the most common, corresponding to a root plus a
single inflectional suffix—the canonical agglutinative pattern.
Figure~\ref{fig:frag_comp} shows the full density comparison.}
\label{tab:frag_dist}
\end{table*}

\begin{figure*}[t]
  \centering
  \includegraphics[width=\linewidth]{figures/fragmentation_comparison.pdf}
  \caption{\textbf{Left}: subword fragment density for Turkish (blue) and
           English (orange). Turkish is right-shifted, reflecting
           morphological complexity.
           \textbf{Right}: mean fragments $\pm$ SD per task.
           The Turkish QA and summarisation tasks both fragment at
           $\approx$2.16 frags/word, confirming that this is a language
           property rather than a task artefact.}
  \label{fig:frag_comp}
\end{figure*}

% ============================================================
\section{Discussion}
\label{sec:discussion}

\subsection{Why Does Turkish Not Lag Behind English?}
\label{sec:disc-turkey}

The central finding—that Turkish acceptance rates ($\alpha = 0.768$)
\emph{exceed} English ($\alpha = 0.716$) despite $1.40\!\times$ higher
subword fragmentation—requires explanation.
We offer a three-part account.

\paragraph{The alignment hypothesis.}
At $T=0.0$, acceptance reduces to: does the draft's greedy prediction
match the target's greedy prediction?
A morphologically complex token that \emph{both} models agree on is
accepted with probability 1.
The key question is whether the draft's top-1 prediction matches the
target's top-1 prediction \emph{more} often for Turkish than for English.
Our results show it does ($76.8\%$ vs.\ $71.6\%$), which we attribute to
the co-training alignment described below.

\paragraph{Same-corpus co-training produces superior argmax alignment.}
The \texttt{ytu-ce-cosmos} draft (117M) and target (774M) are both
pre-trained from scratch on the same 35\,GB Turkish corpus with the same
BPE tokenizer and identical data preprocessing.
Their greedy predictions are therefore strongly correlated: when the
large model assigns highest probability to token $t$, the small model--
trained on identical data--is likely to also rank $t$ first.
The OpenAI GPT-2 small and large models, while sharing architecture and
tokenizer, differ in training data mixture and optimisation details,
leading to lower argmax agreement ($71.6\%$).
This same-corpus alignment advantage more than offsets the modest
additional fragmentation of Turkish morphology.

\paragraph{Evidence from error analysis.}
Section~\ref{sec:error} shows that rejection rates for common function
words and demonstratives (\textit{bu}: 24.1\%, \texttt{-}: 22.4\%,
\textit{ve}: 20.4\%) are elevated compared to punctuation (period: 8.5\%),
consistent with longer-range context dependence.
However, even these elevated rates are modest relative to the overall
acceptance rate of 76.8\%, and the position analysis
(Section~\ref{sec:position}) shows no accumulation of errors over the
generated sequence.
Together, these observations support the interpretation that domain
co-training—not merely shared vocabulary—is the key enabler.

\paragraph{Practical implication.}
Morphological complexity is not a barrier to speculative decoding;
same-corpus co-training can make it an \emph{advantage}.
A Turkish instruction-tuned model paired with a language-matched,
same-corpus draft will outperform an English-paired draft at the same
size ratio.
Our Llama ($\alpha=0.448$) and Qwen ($\alpha=0.370$) supplementary
experiments confirm this: cross-lingual draft--target pairs (Turkish
text, English-base drafts) achieve substantially lower acceptance rates,
supporting the co-training hypothesis.

\subsection{Why Does Turkish Achieve Higher Speedup Than English?}
\label{sec:disc-speedup}

The observed speedup advantage of Turkish over English ($1.250\!\times$
vs.\ $1.159\!\times$) directly reflects the $5.2$\,pp acceptance rate
advantage ($0.768$ vs.\ $0.716$).
Higher acceptance means more tokens are accepted per target call,
amortising the fixed verification cost over a larger yield.

\textit{Same-corpus co-training effect.}
The \texttt{ytu-ce-cosmos} draft (117\,M) and target (774\,M) were
trained on the \emph{identical} 35\,GB Turkish corpus with the same BPE
tokenizer.
Their greedy predictions are therefore more strongly aligned than the
OpenAI GPT-2 small/large pair, which were trained in separate runs on
different data mixes.
At $T=0.0$, acceptance reduces to checking whether draft and target agree
on their argmax; same-corpus training directly increases this agreement.

\textit{Why medium drafts reverse the gain.}
Medium draft models (354\,M) raise acceptance by $+3.7$\,pp (TR) but
cost ${\approx}3\!\times$ more per token than small drafts.
The extra draft computation outweighs the benefit of fewer target calls,
flipping the speedup to a slowdown ($0.918\!\times$).
The theoretical formula $(1+\gamma\alpha)/(1+\gamma/k)$ with $k=6.6$
(small) vs.\ $k=2.2$ (medium) predicts this: at $\alpha=0.768$,
$\gamma=5$, the formula gives $2.58\!\times$ for small and $1.26\!\times$
for medium in an ideal implementation; the observed ratios are
$1.250\!\times$ and $0.918\!\times$, with the gap attributable to
the Python dispatch overhead that is constant regardless of $k$.

\subsection{Asymmetry of the TR/EN Experiment Design}
\label{sec:disc-symmetry}

The Turkish small-draft experiment is replicated with three seeds
$\{42, 123, 456\}$ ($n=1{,}000$ per seed, QA + summarisation) for a
total of 3,000 samples; the English experiment uses a single seed
(seed 42, $n=500$, SQuAD QA only).
This asymmetry is deliberate: Turkish is our primary language of
interest (addressing RQ1--RQ3), so we invest additional runs to
verify stability; English serves as the controlled comparison baseline.
At $T=0.0$ the accept/reject step is deterministic, so a single-seed
English result is sufficient—the acceptance rate would be identical
across seeds ($\sigma_{\alpha} = 0.000$), as the Turkish three-seed
replication confirms.
Turkish spans two task types (XQuAD-TR QA: $\alpha=0.780$; TR-News
summarisation: $\alpha=0.755$), providing within-language
generalisability evidence absent in the English evaluation.

\subsection{Draft Model Size as the Key Efficiency Lever}
\label{sec:disc-scale}

Section~\ref{sec:scale} demonstrates that the draft-to-target parameter
ratio is a stronger predictor of implementation overhead than language
morphology.
Small-draft pairs ($1\!:\!6.6$) achieve clear speedup ($1.250\!\times$
TR, $1.159\!\times$ EN); medium-draft pairs ($1\!:\!2.2$) flip to
slowdown ($0.918\!\times$, $0.859\!\times$) despite higher acceptance.
Acceptance rates differ by only $3.7$--$5.4$\,pp between scales.
A practical design guideline: keep the draft below ${\approx}15\%$ of
the target's parameter count (small draft: 15.1\%, just above the
threshold; medium draft: 45.7\%, clearly too large).
At the emerging 7B/70B scale, a 1B draft (14.3\% ratio) would sit
in the same favorable regime as our 117M/774M small pair.

\subsection{Implications for Other Morphologically Rich Languages}
\label{sec:disc-implications}

Our results are encouraging for Finnish, Hungarian, Arabic, Georgian,
and other agglutinative or highly inflected languages.
The \emph{necessary} condition is a draft model pre-trained natively
on the target language; the \emph{sufficient} condition is a
draft/target size ratio below $\approx\!1\!:\!5$.
If both are met, our evidence suggests that morphological complexity
does not systematically degrade acceptance rates.

Future work should test this hypothesis for morphologically richer
languages such as Finnish (average word length $>$6 morphemes) or Arabic
(templatic root-and-pattern morphology), where the BPE fragmentation
ratio relative to English may exceed 2$\!\times$ rather than the
$1.40\!\times$ observed for Turkish.

\subsection{Cross-Lingual Draft--Target Pairs: Why Do They Fall Short?}
\label{sec:disc-crosslingual}

The supplementary experiments with Llama-3 and Qwen-2.5 reveal a sharp
dichotomy: same-corpus Turkish pairs achieve $\alpha=0.768$, while
cross-lingual pairs (Turkish prompts, English-base drafts) reach only
$\alpha=0.448$ (Llama-3) and $\alpha=0.370$ (Qwen-2.5)—gaps of
$32.0$ and $39.8$\,pp respectively.
This is a more severe penalty than the TR/EN language difference
($5.2$\,pp) and warrants a separate explanatory account.

\paragraph{Vocabulary mismatch as the primary driver.}
Llama-3 and Qwen-2.5 are trained primarily on English and use
English-optimised BPE vocabularies.
When these models process Turkish text, Turkish morphological suffixes
are split into byte-level or character-level fragments not present in
the English vocabulary, producing high-entropy predictions across many
low-probability continuations.
The draft's top-1 prediction for a Turkish suffix then almost never
matches the target's top-1 prediction—not because the models disagree
semantically, but because they have incommensurable distributions over
a fragmented token space.

\paragraph{The domain/instruction-tuning gap.}
Llama-3 and Qwen-2.5 are instruction-tuned on English-dominant corpora,
while the \texttt{ytu-ce-cosmos} target is a base model trained on a
Turkish-only corpus.
The instruction-tuning shifts the target's distribution toward
response-format tokens (headers, lists, affirmations) that the
Turkish-only draft has not seen, creating an additional source of
argmax disagreement beyond vocabulary fragmentation.

\paragraph{What this implies for Turkish LLM serving.}
As multilingual instruction-tuned models with Turkish-adapted
vocabularies become available (e.g., Turkish-fine-tuned Llama-3 with
extended vocabulary), cross-lingual acceptance rates may approach
same-corpus levels.
Until then, the \texttt{ytu-ce-cosmos} family remains the only
ready-to-deploy same-corpus draft--target pair for Turkish, and
the cross-lingual penalty ($-32$\,pp to $-40$\,pp) should be treated
as a hard constraint when choosing a speculative decoding setup.

\subsection{Tokenization Efficiency and the Fragmentation--Acceptance Trade-off}
\label{sec:disc-tokenization}

A natural prediction from tokenization theory is that higher BPE
fragmentation should \emph{hurt} speculative decoding: more tokens must
be proposed and verified per word, increasing the number of decision
points at which a mismatch can terminate the draft.
Our results challenge this prediction.
Turkish tokens are fragmented at $1.40\!\times$ the English rate
(Section~\ref{sec:fragmentation}), yet Turkish acceptance rates exceed English
by $5.2$\,pp.

The resolution is that fragmentation and acceptance are \emph{not}
in direct opposition when the draft and target share a training corpus.
\citet{kaya2022tokenization} showed that tokenization granularity
strongly modulates downstream task performance for Turkish LLMs;
our results show it also modulates the \emph{difficulty} of the
speculation problem.
A Turkish suffix such as \textit{-lerin} that appears as a single
high-frequency BPE token is predicted by both models with high
confidence (since it follows a small set of nominal stems in
learnable distributional patterns), yielding near-certain acceptance.
A rarer English morphological construction expressed across multiple
tokens may produce divergent top-1 predictions between small and
large models even when both models are trained on comparable data.

The practical takeaway is that BPE fragmentation ratios should
\emph{not} be used as a proxy for speculative decoding difficulty in
morphologically rich languages; per-task acceptance rate measurement
is necessary to assess feasibility.

\subsection{Practical Deployment Recommendations}
\label{sec:disc-deployment}

Aggregating the empirical findings into actionable guidance for
practitioners deploying Turkish (or morphologically rich language) LLMs:

\paragraph{Pair draft and target on the same native corpus.}
Co-trained, same-tokenizer draft--target pairs (e.g., \texttt{ytu-ce-cosmos}
117M + 774M) achieve $\alpha=0.768$, yielding $1.25\!\times$ wall-clock
speedup on a single GPU.
Cross-lingual pairs (Turkish prompt, English-base draft) achieve only
$\alpha\approx0.37$--$0.45$, providing little to no benefit.
Until Turkish-native instruction-tuned models are widely available,
the \texttt{ytu-ce-cosmos} family provides the strongest ready-to-use
draft--target pair.

\paragraph{Keep the draft-to-target ratio below $\sim$1:5.}
Our ablation and scale experiments (Sections~\ref{sec:ablation},
\ref{sec:scale}) both confirm that draft models larger than
$\approx\!15$--$20\%$ of the target incur overhead that eliminates
speedup under Python-level serving.
For a 7B target, a 1B draft (14.3\% ratio) is the recommended ceiling.

\paragraph{Set $\gamma = 5$ as the default speculation width.}
The $\gamma$ ablation (Section~\ref{sec:ablation}) shows that speedup
rises steeply from $\gamma=1$ to $\gamma=5$ and plateaus thereafter.
Values of $\gamma\geq7$ provide negligible additional gain while
increasing memory and dispatch cost.
$\gamma=5$ balances throughput and latency across both QA and
summarisation task types.

\paragraph{Expect consistent gains across tasks at $T=0.0$.}
At greedy temperature the accept/reject step is deterministic, so
acceptance rates are reproducible to 3 significant figures ($\sigma=0.000$
across seeds).
Practitioners can estimate expected speedup from a single short
calibration run without statistical uncertainty.

% ============================================================
\section{Limitations}
\label{sec:limitations}

\begin{itemize}[leftmargin=*]

  \item \textbf{English baseline asymmetry.}
        The Turkish experiment uses three random seeds ($n=3{,}000$
        samples, two task types: QA + summarisation), while the English
        baseline uses a single seed ($n=500$, extractive QA only).
        Although $T=0.0$ determinism makes seed variation irrelevant for
        acceptance rate, the \emph{task} asymmetry remains: we cannot
        directly attribute the TR/EN acceptance gap to language rather
        than to the different text-generation regimes (open-ended
        summarisation vs.\ constrained QA extraction).
        A fully symmetric comparison—English with both task types,
        three seeds, and 3,000 samples—would provide a cleaner
        cross-linguistic contrast and is the most important single
        extension of this work.

  \item \textbf{Morpheme-category acceptance analysis.}
        Our error analysis (Section~\ref{sec:error}) identifies which
        \emph{token strings} are rejected most frequently, but does not
        link rejections to morphological categories
        (ROOT, NOMINAL\_SUFFIX, VERBAL\_SUFFIX, TENSE, AGREEMENT, etc.).
        This is the most theoretically interesting open question: if
        verbal agreement suffixes carry higher rejection rates than
        nominal case markers, that would reveal whether speculative
        decoding difficulty tracks syntactic dependency length rather
        than surface token frequency.
        The obstacle is that GPT-2-generated Turkish text contains
        hallucination and code-switching that prevents reliable
        morphological parsing of the output side.
        A study using \emph{input} prompt tokens—sourced from clean
        Wikipedia or UD treebank text—would yield clean morpheme-level
        breakdowns and should be a priority for future work.

  \item \textbf{Cross-lingual model families not deeply analysed.}
        The supplementary Llama-3 ($\alpha=0.448$) and Qwen-2.5
        ($\alpha=0.370$) experiments establish that cross-lingual
        draft--target pairs perform substantially worse than same-corpus
        pairs ($\alpha=0.768$), but these results are reported without
        detailed analysis.
        It remains unknown whether the gap is driven by
        \emph{vocabulary mismatch} (different BPE vocabularies forcing
        the draft to hallucinate tokens absent from the target's
        distribution), \emph{embedding space misalignment}, or
        \emph{training data domain shift}.
        Each cause implies a different remedy: adapting the tokenizer,
        aligning embedding spaces via distillation, or domain-specific
        fine-tuning.
        A systematic ablation isolating these factors—using Turkish-fine-tuned
        Llama and Qwen variants as they become available—would substantially
        extend the practical reach of this work.

  \item \textbf{Greedy decoding only ($T = 0.0$).}
        All experiments use temperature $T=0.0$, making acceptance
        fully deterministic and ensuring perfect output fidelity.
        In production, however, temperature is rarely zero:
        instruction-following and dialogue systems typically sample at
        $T \in [0.6, 1.0]$ to increase output diversity.
        At $T > 0$ the accept/reject criterion becomes stochastic
        \citep{leviathan2023fast}: a draft token $d_i$ is accepted
        with probability $\min(1, p_{\text{target}}(d_i) /
        p_{\text{draft}}(d_i))$, and any rejection triggers a
        corrective sample from the residual distribution.
        Turkish morphology may interact non-trivially with this
        stochastic criterion: low-probability suffix continuations
        that \emph{both} models assign high uncertainty could produce
        acceptance rates that diverge more strongly across temperature
        levels than the corresponding English constructions.
        Extending the experimental design to $T \in \{0.3, 0.7, 1.0\}$
        is a necessary step before deployment claims can be made for
        production Turkish dialogue systems.

  \item \textbf{Task scope.}
        We evaluate on extractive QA and abstractive summarisation.
        Open-ended dialogue or code generation may produce different
        acceptance rate profiles due to higher output-space entropy.

  \item \textbf{Implementation overhead for medium-draft models.}
        The wall-clock slowdown observed for medium-draft pairs
        ($0.918\!\times$, $0.859\!\times$) reflects both the unfavourable
        size ratio and residual Python dispatch overhead.
        A production C++/CUDA implementation would close this gap;
        the acceptance rate findings are implementation-independent.

  \item \textbf{Single GPU, single hardware platform.}
        All experiments run on a Google Colab NVIDIA T4 (16\,GB VRAM,
        320\,GB/s bandwidth).
        Results--particularly the absolute speedup values--may differ on
        A100/H100 hardware, where memory bandwidth is 6--12$\!\times$
        higher and the draft phase becomes proportionally cheaper.
        The \emph{acceptance rate} findings, being hardware-independent,
        are expected to generalise.

\end{itemize}

% ============================================================
\section{Conclusion}
\label{sec:conclusion}

We have presented the first systematic study of speculative decoding on
Turkish, an agglutinative language with substantially higher subword
fragmentation than English.
Our key findings are:

\begin{enumerate}[leftmargin=*]
  \item Turkish acceptance rate \emph{exceeds} English ($0.768$ vs.\
        $0.716$; $\Delta=5.2$\,pp, $p \approx 0$) despite $1.40\!\times$
        higher subword fragmentation—the paper's primary finding.
  \item At $T=0.0$, acceptance is fully deterministic ($\sigma=0.000$
        across three seeds), guaranteeing perfect reproducibility.
  \item Both small-draft pairs achieve statistically significant speedup:
        $1.250\!\times$ (TR) and $1.159\!\times$ (EN), both
        $p < 10^{-40}$.
  \item Draft model size is the primary efficiency lever: medium drafts
        raise acceptance by $3.7$--$5.4$\,pp but flip to slowdown
        ($0.918\!\times$ TR, $0.859\!\times$ EN) due to heavier
        draft computation.
  \item Acceptance rate increases slightly across position
        ($87.5\%\!\to\!93.3\%\!\to\!93.7\%$ early$\to$mid$\to$late
        thirds), ruling out accumulation of morphological difficulty.
  \item Output is empirically lossless: ROUGE-1 $\Delta < 0.001$
        and BLEU $\Delta < 0.001$ between greedy and speculative.
\end{enumerate}

These findings establish that speculative decoding is fully viable for
agglutinative Turkish NLP, achieving $1.25\!\times$ speedup with a
same-tokenizer, natively pre-trained draft--target pair.
At $T=0.0$, acceptance is deterministic and output is identical to
greedy decoding.
Error analysis reveals that demonstratives and function words carry the
highest rejection rates (24.1\%), consistent with their variable
morphosyntactic contexts.

Four extensions stand out.
A fully symmetric English baseline (both task types, $n=3{,}000$) would
disentangle language effects from task-regime differences.
Morpheme-category acceptance analysis on clean UD treebank prompts would
answer which morphological categories drive rejection—the paper's most
open theoretical question.
Experiments at $T > 0$ are necessary before claims extend to production
dialogue systems, where greedy decoding is rarely used.
Finally, Turkish-native Llama and Qwen variants would allow ablation of
vocabulary mismatch vs.\ domain shift in cross-lingual pairs, and
extension to Finnish, Hungarian, and Arabic would test the generality of
the co-training alignment hypothesis.

% ============================================================
\section*{Acknowledgements}
The author thanks the \texttt{ytu-ce-cosmos} team at Y{\i}ld{\i}z
Technical University for releasing the Turkish GPT-2 model family under
an open licence, and Google for providing GPU compute through the Colab
platform.

\end{multicols}

% ── References (single column) ───────────────────────────────
\bibliography{references}

\end{document}
